\chapter{Discussion}\label{chapter:chap8}

\section{Interpretation of Results}\label{sect:interpretation}

The key findings from the analysis is the following:

8.4\% of analyzed crates contain at least one vulnerability, with an average of 4.37 per vulnerable crate. This is lower than the npm rate, however it remains a concern

The top 10 dependent crates have over 350.000 dependent crates.

Over 393.059 secrets were detected by gitleaks, however the majority of them are false positives of generic API key patterns, and probably represent test and constants. The non-generic detections are a small number and should require manual analysis.

Around 4.6\% of crates contain build scripts, with over 240 performing network operations.

\section{Comparison with Existing Approaches}\label{sect:comparison}

TODO: After rerun of framework

\section{Limitations}\label{sect:limitations}

\subsection{Detection Limitations}\label{subsect:detection-limitations}

False negatives:

- Novel patterns and attacks can evade our heuristics

- Custom encoding might be done to try to evade obfuscation detection

- Social engineering and account overtaking are not addressed

False positives:

- Legitimate crates with unusual behavior that still remain benign

- Test code may trigger alerts

- High entryopy with legitimate usecases

- Common naming patterns can produce false alerts of the typosquatting features

\subsection{LLM-Specific Limitations}\label{subsect:llm-limitations}

- LLMs might hallucinate vulnerabilities or hallucinate that benign code is malicious or viceversa

- The very limited context window is a problem due to some libraries having large codebases, which are chunked on the current implementation, which can make the model miss cross-file malicious behavior (behavior that present in two separate files is benign on it's own but either imports or api usage makes it pontentially malicious)

- Reproducibility is another limitation since LLM outputs vary even with low or zero temperature parameter.

- Prompt injection is another problem, adversaries might contain specially encoded unicode characters or try to do prompt injection in comments to evade the model

\subsection{Scalability and Resource Limitations}\label{subsect:scalability-limitations}

LLm inference is extremely computationally expensive for an ecosystem deployment.

Analysis time can take a very long time, especially on large crates, adversaries could do a DoW attack by spamming fake crates with large codebases to attack crates.io.

Storage is another problem if storing crates is required, besides us only cloning with a depth of one and only one branch, some repositories are still large

Besides storage, network bandwidth is also a potential problem, especially if the git services used rate limit us.

\subsection{Ground Truth Limitations}\label{subsect:ground-truth-limitations}

Validation is extremely hard to do, since only a very small number of existing crates have been documented in a public way as malicious

We also have an unknown rate for false negatives since we cannot know how many attacks exists undetected

\subsection{Generalizability Limitations}\label{subsect:generalizability-limitations}

The framework is designed for crates.io and Rust, adoption for other ecosystems would require a full rewrite.

Besides the tools like cargo audit being adapted for Rust, also the build.rs heuristics are also language specific

\section{Threats to Validity}\label{sect:threats}

The CVSS scoring and string similarity algorithms have been tested via unit tests, however we can never know if there are any edge cases

Results are limited to crates with valid repository links in the metadata, some crates have been moved from the link from the metadata and moved to docs.rs (i.e. Find\_Matt\_C)

The LLM score is only a guide and cannot guarantee anything, be it either the presence of malivious behavior or its absence.